\documentclass[a4paper,12pt,fleqn]{scrartcl}%fleqn pour aligner les equations à gauche
\usepackage{../../Styles/Cours}


\newcommand{\chnum}{Chapitre 8}
\newcommand{\chtitle}{Réaction chimiques forcées}
\newcommand{\dtype}{Cours}


\begin{document}
	\maketitle
\section{Électrolyse}
\subsection{Principe}
Une réaction d'oxydorection spontanée a lieu dans un sens bien déterminé entre l'oxydant d'un couple et le réducteur d'un autre couple. Le système chimique se rapprochant de son état d'équilibre, cette réaction ne peut pas s'inverser naturellement.\\
Si un générateur de courant continu est branché correctement, il peut inverser les transformations ayant lieu dans le système, permettant ainsi la conversion d'énergie électrique en énergie chimique. Cette technique est appelée électrolyse.

\subsection{Transformation chimique forcée}
Pour inverser une réaction d'oxydoreduction, il faut que le transfert d'électron entre les couples change de sens. Cette inversion du sens de la réaction est rendue possible grâce à un générateur.\\

\emph{Exemple :}\\
Pour les couples \ce{I_2(aq)}/\ce{I-(aq)} et \ce{Zn^{2+}(aq)}/\ce{Zn(s)} :\\
\textbf{Sens de la réaction spontanée :} \ce{Zn(s) + I_2(aq) -> 2I-(aq) + Zn^{2+}(aq)}\\
\textbf{Sens de la réaction forcée :} \ce{2I-(aq) + Zn^{2+}(aq) -> Zn(s) + I_2(aq)}\\
Ce qui implique les réactions aux électrodes :\\
\ce{2I-(aq) -> I_2(aq) + 2e-} (sens de l'oxydation)\\
\ce{Zn^{2+}(aq) + 2e- -> Zn(s)}\\

\vspace{1em}
L'équation de la réaction réalisée lors d'une électrolyse est obtenue par la combinaison des demi-équations de réactions d'oxydation et de réduction faites aux électrodes.\\
Lors d'une réaction spontanée, la proportion des produits augmente au cours de l'évolution, donc le quotient de réaction $Q_r$ augmente et tend ver la constante d'équilibre $K$\\
Dans le cas d'une réaction forcée dans le sens inverse, selon ce même quotient $Q_r$, la proportion des produits diminue, donc $Q_r$ diminue et s'éloigne de $K$.
\end{document}